\documentclass[9.5pt, a4paper]{article}

% ── PACKAGES ──────────────────────────────────────────────────
\usepackage[utf8]{inputenc}
\usepackage[T1]{fontenc}
\usepackage{geometry}
\geometry{
  top=1.0cm,
  bottom=1.0cm,
  left=1.4cm,
  right=1.4cm
}
\usepackage{hyperref}
\hypersetup{
  colorlinks=true,
  urlcolor=blue,
  linkcolor=blue,
  pdftitle={LECA Resurrection Protocol — One Page Summary},
  pdfauthor={Eric Robert Lawson / OrganismCore}
}
\usepackage{microtype}
\usepackage{parskip}
\usepackage{booktabs}
\usepackage{array}
\usepackage{tabularx}
\usepackage{xcolor}
\usepackage{mdframed}
\usepackage{multicol}
\usepackage{enumitem}
\usepackage{titlesec}
\usepackage{lmodern}

% ── COLOURS ───────────────────────────────────────────────────
\definecolor{headerblue}{RGB}{20,60,140}
\definecolor{statementbox}{RGB}{20,60,140}
\definecolor{rulecolor}{RGB}{20,60,140}

% ── SECTION FORMAT ────────────────────────────────────────────
\titleformat{\section}
  {\small\bfseries\color{headerblue}}
  {}{0em}{}[\vspace{-0.25em}
  \textcolor{rulecolor}{\rule{\linewidth}{0.4pt}}]
\titlespacing{\section}{0pt}{0.45em}{0.2em}

\setlength{\parskip}{0.2em}
\setlength{\parindent}{0pt}
\setlength{\columnsep}{0.9em}
\pagestyle{empty}

% ══════════════════════════════════════════════════════════════
\begin{document}

% ── HEADER ────────────────────────────────────────────────────
\begin{center}
  {\large\bfseries\color{headerblue}
  THE LECA RESURRECTION PROTOCOL}\\[0.15em]
  {\small\color{headerblue}
  Recovery of the Last Eukaryotic Common Ancestor
  from Living Developmental Programs — One-Page Summary}\\[0.1em]
  {\footnotesize
  Eric Robert Lawson\ $\cdot$\
  OrganismCore (Independent Research)\ $\cdot$\
  ORCID:
  \href{https://orcid.org/0009-0002-0414-6544}
  {0009-0002-0414-6544}\ $\cdot$\
  \href{mailto:OrganismCore@proton.me}
  {OrganismCore@proton.me}\ $\cdot$\
  2026-03-12}
\end{center}

\vspace{0.15em}

% ── THE CLAIM ─────────────────────────────────────────────────
\begin{mdframed}[
  backgroundcolor=statementbox,
  linecolor=statementbox,
  linewidth=0pt,
  innertopmargin=0.35em,
  innerbottommargin=0.35em,
  innerleftmargin=0.7em,
  innerrightmargin=0.7em
]
\begin{center}
\footnotesize\bfseries\color{white}
Every eukaryotic organism carries the Last Eukaryotic Common
Ancestor (\textasciitilde{}1.5\,Ga) as the base state of its developmental
program. Arresting that program before the first multicellular
step produces a LECA-grade organism. The same result from
an animal, a plant, and a fungus is morphologically identical.
Pre-registered 2026-03-12 with \textbf{zero biological
contradictions}. The experiment costs \textbf{\$200--400} and
takes \textbf{days}.
\end{center}
\end{mdframed}

\vspace{0.3em}

% ── TWO COLUMNS ──────────────────��────────────────────────────
\begin{multicols}{2}

% ── LEFT COLUMN ───────────────────────────────────────────────

\section{The Theorem}

The developmental trajectory from zygote to
adult IS the evolutionary trajectory of the
lineage, compressed into developmental time
(the Reproduction Theorem). Therefore:

\smallskip
{\small
\begin{tabular}{p{0.55cm}p{5.5cm}}
\textbf{E} &
\textbf{Evolutionary grade} — the attractor the
lineage occupied at that point in its history. \\[0.2em]
\textbf{D} &
\textbf{Developmental position} — the equivalent
stage in the developmental program today. \\
\end{tabular}
}

\smallskip
{\small
$E \equiv D$ (compressed). Arresting at position $D$
stabilises grade $E$. The LECA sits at $D = 0$:
before the first multicellular step, in every
eukaryote alive.
\textbf{Arrest gives access. The information was never lost.}
}

{\footnotesize
Pre-registration DOI:\
\href{https://doi.org/10.5281/zenodo.18986790}
{\texttt{10.5281/zenodo.18986790}}
}

\section{The Three Pre-Registered DOIs}

{\footnotesize
\begin{tabular}{p{3.6cm}p{2.4cm}}
\toprule
\textbf{Document} & \textbf{DOI} \\
\midrule
Pre-Registration v1.0 &
\href{https://doi.org/10.5281/zenodo.18986790}
{\texttt{...18986790}} \\[0.15em]
Registered Amendment A1 &
\href{https://doi.org/10.5281/zenodo.18989573}
{\texttt{...18989573}} \\[0.15em]
The Ribosome Is the LUCA &
\href{https://doi.org/10.5281/zenodo.18988844}
{\texttt{...18988844}} \\
\bottomrule
\end{tabular}

\smallskip
Repository:\
\href{https://github.com/Eric-Robert-Lawson/attractor-oncology}
{\texttt{github.com/Eric-Robert-Lawson/}
\texttt{attractor-oncology}}
}

\section{The Five Pre-Registered Predictions}

{\small
\begin{tabular}{p{0.35cm}p{5.7cm}}
\textbf{P1} &
Developmental arrest of any eukaryote before
multicellular organisation produces a single-cell,
nucleated, motile-capable organism. \\[0.2em]
\textbf{P2} &
Arrested organisms from animal, plant, and fungal
starting material are \textit{morphologically
indistinguishable} under light microscopy. \\[0.2em]
\textbf{P3} &
They are \textit{genomically distinct} by sequencing
— the same attractor, different substrates. \\[0.2em]
\textbf{P4} &
Transcriptomics: LECA-grade conserved genes
active; all lineage-specific post-LECA genes
epigenetically silenced. \\[0.2em]
\textbf{P5} &
The blind panel cannot identify starting kingdom
from morphology alone (above 1/3 chance). \\
\end{tabular}
}

\columnbreak

% ── RIGHT COLUMN ──────────────────────────────────────────────

\section{The Three Experimental Arms}

{\small
\begin{tabular}{p{1.1cm}p{2.05cm}p{1.0cm}p{1.35cm}}
\toprule
\textbf{Arm} & \textbf{Start} &
\textbf{Cost} & \textbf{Timeline} \\
\midrule
A — Fungal &
\textit{S.~cerevisiae} spores &
\$200--350 & 24--72\,h \\[0.15em]
B — Animal &
Mouse zygote &
\$400--600 & 24--72\,h \\[0.15em]
C — Plant &
\textit{Arabidopsis} zygote &
\$200--400 & 48--72\,h \\
\midrule
\multicolumn{4}{p{5.5cm}}{
\footnotesize
Full cross-kingdom run (A+B+C): \$800--1,350.
Arm~A alone confirms P1. All three confirm P2--P5.
\textbf{Arm~A is executable at any BSL-1
community biology laboratory today.}}\\
\bottomrule
\end{tabular}
}

\section{Arm A — What You Actually Do}

{\small
\textbf{Day 0:} Sporulate \textit{S.~cerevisiae}
in 1\% KAc medium, 25\,°C, 48--72\,h.
Confirm asci by microscopy.

\textbf{Day 3:} Wash spores into arrest medium
(10\,mM NaCl, 5\,mM KCl, 2\,mM MgCl$_2$,
1\,mM CaCl$_2$, 0.1\,mM FeSO$_4$,
1\,mM KH$_2$PO$_4$, 0.1\% glucose,
\textbf{no nitrogen source},
pH\,6.5--7.0). Culture at 20--22\,°C,
$<$3\% O$_2$.

\textbf{Days 4--10:} Observe at 24\,h, 48\,h,
72\,h. Photograph by phase contrast (40$\times$).
DAPI stain to confirm nuclear integrity.
Check: single-cell form maintained,
no germ tube, viable ($\geq$80\%).

\textbf{Pass criterion:}
$\geq$80\% single-cell, nucleus confirmed,
$\leq$10\% germination, viable at 72\,h.

\textbf{Equipment:} Laminar flow cabinet,
incubator, compound microscope (40$\times$
phase contrast), fluorescence microscope
for DAPI. All standard BSL-1.
}

\section{What the Experiment Establishes}

{\small
\begin{tabular}{p{2.7cm}p{2.7cm}}
\textbf{One arm confirms} & \textbf{Three arms confirm} \\
\midrule
Arrest accesses LECA grade & Cross-kingdom convergence \\
Protocol is executable & Attractor is real \\
P1 confirmed & P1--P5 confirmed \\
Pre-reg validated (P1) & Pre-reg fully validated \\
Publishable result & Historic result \\
\end{tabular}
}

\smallskip
{\small
The oldest recoverable organism is the easiest to
recover: the LECA sits at developmental position
zero. \textbf{No ancient DNA. No special equipment.
No institutional affiliation required.}
The ribosome's peptidyl transferase centre — the
LUCA's continuously executing Row~3 identity
(DOI \texttt{...18988844}) — will be present
in every organism this protocol recovers,
regardless of starting kingdom.
}

\end{multicols}

\vspace{0.1em}
\noindent\textcolor{rulecolor}
{\rule{\linewidth}{0.4pt}}

% ── FOOTER ────────────────────────────────────────────────────
\begin{center}
\footnotesize
\textit{The geometry came first.\ \
The developmental program never lost it.\ \
The experiment stops the clock and looks.}\\[0.2em]
\href{https://github.com/Eric-Robert-Lawson/attractor-oncology}
{\texttt{github.com/Eric-Robert-Lawson/attractor-oncology}}
\ $\cdot$\
\href{https://doi.org/10.5281/zenodo.18986790}
{\texttt{doi.org/10.5281/zenodo.18986790}}
\ $\cdot$\
\href{https://orcid.org/0009-0002-0414-6544}
{\texttt{ORCID: 0009-0002-0414-6544}}
\end{center}

\end{document}
